\label{sec:experiments}

\subsection{Dataset}

Daily BTC/USD OHLCV data sourced from Kaggle spanning 2014 to 2022 was used.
The training set covers the pre-2020 period (approximately 2014--2019) and
the test set covers 2020--2022. The split was not chosen arbitrarily:
post-2020 data encompasses the COVID-19 market crash of March 2020, the 2021
bull run, and the 2022 bear market, constituting three qualitatively distinct
market regimes. This is an intentional out-of-distribution stress test rather
than a random holdout. Its purpose is to measure whether parameter
configurations optimised on one regime generalise when the underlying market
dynamics shift substantially --- a more meaningful evaluation of strategy
robustness than a matched-regime random split would provide.

\subsection{Experimental Design}

The full factorial design comprised 6~algorithms $\times$ 5~strategies
$\times$ 2~initialisation modes $\times$ 30~runs = 1800~total trials. Two
initialisation modes were tested: fixed initialisation, in which all runs
begin from a deterministic seed position, and random initialisation, in which
each run begins from a position drawn uniformly at random from the search
space. The paired comparison of fixed versus random initialisation tests
sensitivity to starting conditions. Strong dependence on initialisation would
indicate algorithmic fragility, with convergence driven by the seed rather
than by the search mechanism. A run count of 30 was chosen as a practical
compromise between statistical power and the total computational cost of 1800
trials; it is sufficient to apply non-parametric distributional tests but
modest relative to what would be preferred for high-variance optimisers.

\subsection{Statistical Analysis}

The Kruskal-Wallis $H$-test was used in preference to one-way ANOVA because
the structural properties of optimiser output systematically violate the
normality assumption required by ANOVA. Optimisers can produce runs that
converge to the degenerate no-trade optimum (fitness near \$1000, the
starting balance), occasional jackpot runs that far exceed the typical
distribution, and heavy-tailed spread in between. These outliers are not
noise --- they reflect real behavioural modes of the algorithms --- and would
distort ANOVA $F$-statistics. Kruskal-Wallis operates on ranks and is robust
to such distributional shapes.

For each of the ten strategy $\times$ initialisation mode combinations, the
Kruskal-Wallis test was applied with algorithm as the factor variable,
holding strategy and initialisation constant. This isolates the algorithm
effect from strategy-level confounds: testing across all strategies
simultaneously would conflate the variance introduced by different fitness
landscape shapes with the variance attributable to algorithmic differences.
Pairwise post-hoc comparisons used Dunn's test with Bonferroni correction. A
significance threshold of $\alpha = 0.05$ was applied throughout.
